% Презентација за одбрану мастер рада.
% Заснована на finbeamer класи (github.com/elektrotehnika/finlatex),
% коју је ментор препоручио 1. марта 2026.
%
% Компилација (latexmk не ради без Perl-а):
%     pdflatex -interaction=nonstopmode presentation.tex
%     pdflatex -interaction=nonstopmode presentation.tex
%
\documentclass[aspectratio=169]{beamer}

\usepackage[T1, T2A]{fontenc}
\usepackage[english, serbianc]{babel}
\usepackage{tikz}
\usepackage{booktabs}
\usepackage{graphicx}

% ── Боје Универзитета ────────────────────────────────────────────────────────
\definecolor{UniBurgundy}{HTML}{970B30}
\definecolor{UniBurgundyDark}{HTML}{5C071E}
\definecolor{UniRose}{HTML}{F4E9EC}
\definecolor{UniInk}{HTML}{241F20}
\definecolor{UniMuted}{HTML}{6D6265}

% A clean university treatment: the burgundy of the official mark is the only
% structural accent, while charts retain their own analytical colour scales.
\usetheme{default}
\setbeamercolor{normal text}{fg=UniInk,bg=white}
\setbeamercolor{structure}{fg=UniBurgundy}
\setbeamercolor{frametitle}{fg=white,bg=UniBurgundy}
\setbeamerfont{frametitle}{series=\bfseries,size=\large}
\setbeamercolor{title}{fg=UniBurgundy}
\setbeamercolor{author}{fg=UniInk}
\setbeamercolor{date}{fg=UniMuted}
\setbeamercolor{section in toc}{fg=UniBurgundy}
\setbeamercolor{subsection in toc}{fg=UniMuted}
\setbeamercolor{itemize item}{fg=UniBurgundy}
\setbeamercolor{itemize subitem}{fg=UniBurgundy}
\setbeamercolor{enumerate item}{fg=UniBurgundy}
\setbeamercolor{block title}{fg=white,bg=UniBurgundy}
\setbeamercolor{block body}{fg=UniInk,bg=UniRose}
\setbeamertemplate{headline}{}
\setbeamertemplate{navigation symbols}{}
\setbeamertemplate{caption}{\insertcaption}
\setbeamertemplate{frametitle}{\nointerlineskip\hbox{\begin{beamercolorbox}[wd=\paperwidth,ht=0.82cm,dp=0.16cm,leftskip=0.55cm]{frametitle}\usebeamerfont{frametitle}\insertframetitle\end{beamercolorbox}}}

% ── Нумерација слајдова ──────────────────────────────────────────────────────
\newcommand\numbered{\setbeamertemplate{footline}{\leavevmode\hbox{\begin{beamercolorbox}[wd=\paperwidth,ht=.34cm,dp=.12cm,rightskip=.38cm]{author in head/foot}\color{UniBurgundy}\scriptsize\hfill\insertframenumber\,/\,\inserttotalframenumber\end{beamercolorbox}}}}
\newcommand\samenumbered{\numbered\addtocounter{framenumber}{-1}}
\newcommand\unnumbered{\setbeamertemplate{footline}{}\addtocounter{framenumber}{-1}}

% ── Лого ─────────────────────────────────────────────────────────────────────
% Користи се НОВИ лого Универзитета (исти као у раду), а не стари грб из
% шаблона — студентска служба је 11. јуна изричито упозорила да се рад не сме
% штампати са старим грбом.
\newcommand{\titlelogouni}{\def\titlelogo{\begin{figure}
\includegraphics[height=52pt]{./slike/grb/unikg_logo}\end{figure}}}
\titlelogouni

\newcommand{\footlogo}[1]{\setbeamertemplate{background canvas}{%
\ifnum\insertframenumber>0
    \begin{tikzpicture}[overlay, remember picture]
        \node[anchor=south west] at (current page.south west) {\includegraphics[width=20pt]{#1}};
    \end{tikzpicture}%
\fi}}
\footlogo{slike/grb/unikg_logo.png}

% ── Подаци о раду ────────────────────────────────────────────────────────────
\title{\large Предвиђање временски променљивих међутржишних зависности\\између криптовалута и традиционалне имовине\\применом машинског учења}
\author{Богдан Бабаев\texorpdfstring{\vspace{-1.5em}}{}}
\institute{\textbf{Мастер академске студије --- Вештачка интелигенција}\\Универзитет у Крагујевцу%
\vspace{4pt}\titlelogo\vspace{-12pt}}
\date{Крагујевац, 10. септембар 2026.}

\newcommand{\FIG}{../outputs/figures}

\begin{document}
\numbered

{\unnumbered
\begin{frame}[plain]
    \centering
    \vspace{0.45cm}
    {\color{UniBurgundy}\large Предвиђање временски променљивих међутржишних зависности\\
    између криптовалута и традиционалне имовине\\
    применом машинског учења\par}
    \vspace{0.55cm}
    {\large Богдан Бабаев\par}
    \vspace{0.12cm}
    {\small\textbf{Мастер академске студије --- Вештачка интелигенција}\\
    Универзитет у Крагујевцу\par}
    \vspace{0.18cm}
    {\small Ментор: ред. проф. др Владимир М. Миловановић\par}
    \vspace{0.22cm}
    
\includegraphics[height=52pt]{./slike/grb/unikg_logo.png}\par
    \vfill
    {\color{UniMuted}\small Крагујевац, 10. септембар 2026.}
\end{frame}}

\section*{Садржај}
\begin{frame}{\secname}
    \begin{columns}[c]
        \begin{column}{0.72\linewidth}
            \tableofcontents[hideallsubsections]
        \end{column}
        \begin{column}{0.28\linewidth}
            \centering
            
\includegraphics[width=0.78\linewidth]{slike/qr/master_thesis_site_qr.png}\\[4pt]
            {\scriptsize Интерактивна презентација рада}
        \end{column}
    \end{columns}
\end{frame}

% ─────────────────────────────────────────────────────────────────────────────
\section{Проблем}

\begin{frame}{Зашто прогнозирати зависност}
    \begin{itemize}
        \item Корелација је основа диверсификације, хеџинга и портфолио-ризика, док је анализа средње вредности и варијансе често третира као константу
        \item Она није константна. Корелације асиметрично реагују на лоше вести и расту када тржишта падају, баш када је диверсификација најпотребнија
        \item Bitcoin додатно наглашава проблем: њиме се тргује глобално и непрекидно, а његова улога на спектру ризика није фиксна
        \item Зато је предмет истраживања \textbf{променљива} величина. Једна корелација за цео узорак сакрила би важне тржишне режиме
    \end{itemize}
\end{frame}

\begin{frame}{Истраживачка питања}
    \begin{enumerate}
        \item Да ли је склизна зависност прогнозива ван узорка?
        \item Да ли машинско учење надмашује DCC-GARCH модел без цурења будућих информација?
        \item Да ли надмашује једноставне моделе постојаности, као што је AR(1)?
        \item Да ли прогнозе могу бити основа за сигнал упозорења на стрес за инвеститора?
    \end{enumerate}
\end{frame}

% ─────────────────────────────────────────────────────────────────────────────
\section{Подаци и циљна променљива}

\begin{frame}{Подаци}
    \begin{columns}
        \begin{column}{0.46\linewidth}
            \begin{itemize}
                \item Основни актив: BTC-USD
                \item Упарен са \^{}GSPC, \^{}IXIC, GLD, SLV, UUP и ETH-USD
                \item Од 9. новембра 2017. до 16. маја 2026, 3\,053 дневна посматрања
                \item Почетни датум одређује Ethereum, јер раније нема цена ETH-USD
                \item Дневни лог-приноси, заједнички временски индекс и без празнина после поравнања
            \end{itemize}
        \end{column}
        \begin{column}{0.54\linewidth}
            \includegraphics[width=\linewidth]{\FIG/eda_normalized_prices.png}
        \end{column}
    \end{columns}
\end{frame}

\begin{frame}{Циљна променљива: склизна зависност и Fisher-$z$}
    \begin{columns}
        \begin{column}{0.44\linewidth}
            \begin{itemize}
                \item Склизна Pearson-ова корелација у прозорима $w = 14, 30, 60, 90$ дана
                \item Fisher-$z$ пресликава $(-1,1)$ на реалну праву, па се RMSE коректно понаша близу граница
                \item Хоризонт прогнозе $h = 1$
                \item 6 парова $\times$ 4 прозора $=$ \textbf{24 експеримента}
            \end{itemize}
        \end{column}
        \begin{column}{0.56\linewidth}
            \includegraphics[width=\linewidth]{\FIG/eda_fisher_transform.png}
            \medskip
            {\scriptsize Ограничена корелација је лево, а Fisher-$z$ трансформација је десно.}
        \end{column}
    \end{columns}
\end{frame}

% ─────────────────────────────────────────────────────────────────────────────
\section{Методологија}

\begin{frame}{\secname}
    \begin{description}
        \item[Walk-forward] проширујући прозор, $\mathrm{min\_train} = 800$, поновно оцењивање на сваких 20 корака
        \item[Модели] Naive, AR(1), HAR, Ridge, ElasticNet, RF, GBM, XGBoost и Ensemble
        \item[Референтни модел] DCC-GARCH(1,1), поново оцењен унутар петље
        \item[Тестирање] Diebold--Mariano тест на упареним разликама губитка
    \end{description}
    \bigskip
    \begin{itemize}
        \item Сваки корак који би могао открити будућност --- конструкција циља, поравнање обележја, оцена референтног модела и генерисање сигнала --- извршава се \emph{унутар} петље
        \item Хиперпараметри су једном изабрани временски унакрсном валидацијом за BTC--S\&P\,500, $w=30$; тај репрезентативни случај је зато истраживачки, а не независна потврдна оцена
    \end{itemize}
\end{frame}

\begin{frame}{Репрезентативна прогноза}
    \begin{center}
        \includegraphics[width=0.70\linewidth]{\FIG/thesis_forecast_corr_BTC-USD_^GSPC_w30_fisher_z.png}
    \end{center}
    \vspace{-4pt}
    {\scriptsize Bitcoin--S\&P\,500, $w = 30$. Модели прате путању, а грешке се групишу при променама режима.}
\end{frame}

% ─────────────────────────────────────────────────────────────────────────────
\section{Резултати}

\begin{frame}{Рангирање модела}
    \small
    \begin{center}
    \begin{tabular}{lrrr}
        \toprule
        Модел & RMSE & MAE & $R^2$ \\
        \midrule
        Ridge              & \textbf{0.0656} & 0.0419 & 0.9432 \\
        AR(1)              & 0.0659 & 0.0403 & 0.9424 \\
        HAR                & 0.0659 & 0.0405 & 0.9424 \\
        Naive (последња вредност) & 0.0666 & \textbf{0.0395} & 0.9412 \\
        \midrule
        GBM                & 0.0886 & 0.0625 & 0.8983 \\
        Random Forest      & 0.0905 & 0.0636 & 0.8936 \\
        XGBoost            & 0.1241 & 0.0864 & 0.7923 \\
        \midrule
        DCC-GARCH          & 0.2136 & 0.1683 & 0.3715 \\
        \bottomrule
    \end{tabular}
    \end{center}
    \medskip
    {\small Разлика од 0.0003 на врху није значајна: сваки bootstrap интервал има ширину око 0.009, а по MAE метрици Naive је први. Табела подржава \textbf{групе модела}, а не редослед унутар најбоље групе.}
\end{frame}

\begin{frame}{Зашто постојаност побеђује}
    \begin{columns}
        \begin{column}{0.55\linewidth}
            \begin{itemize}
                \item Склизна корелација од 30 дана дели 29 од 30 посматрања са јучерашњом вредношћу
                \item После Fisher-$z$ трансформације први лаг већ обухвата већину прогнозиве варијансе
                \item Упарени DM, Ridge наспрам Naive: \textbf{6} победа, \textbf{9} пораза и \textbf{9} нерешених од 24 случаја
                \item GBM, Random Forest и XGBoost имају већи RMSE од Naive модела у свих 24 експеримента
            \end{itemize}
        \end{column}
        \begin{column}{0.45\linewidth}
            \includegraphics[width=\linewidth]{\FIG/model_improvement_over_naive.png}
        \end{column}
    \end{columns}
    \medskip
    {\small Зато $R^2 = 0.94$ говори више о циљној променљивој него о моделу.}
\end{frame}

\begin{frame}{Аблација група обележја}
    \scriptsize
    \begin{center}
    \begin{tabular}{lrrrr}
        \toprule
        Скуп обележја & \# & Ridge & Random Forest & XGBoost \\
        \midrule
        G1 само лаг зависности        &  1 & 0.0899 & 0.0986 & \textbf{0.0980} \\
        G2 $+$ историја зависности    & 11 & \textbf{0.0649} & \textbf{0.0884} & 0.1030 \\
        G3 $+$ приноси Bitcoin-а      & 20 & 0.0649 & 0.0891 & 0.1058 \\
        G4 $+$ приноси повезаних актива & 29 & 0.0650 & 0.0902 & 0.1078 \\
        G5 пуни модел                 & 32 & 0.0651 & 0.0910 & 0.1134 \\
        \bottomrule
    \end{tabular}
    \end{center}
    \normalsize
    \medskip
    \begin{itemize}
        \item Једна колона, без међуактивних информација, већ даје Ridge $R^2 = 0.885$
        \item Помаже само сопствена историја циља; од G2 надаље две трећине матрице обележја не додаје ништа
        \item Код стабала сваки додати блок \emph{штети}: XGBoost се погоршава за 16\% од G1 до G5
    \end{itemize}
\end{frame}

\begin{frame}{Машинско учење и DCC-GARCH}
    \begin{columns}
        \begin{column}{0.52\linewidth}
            \includegraphics[width=\linewidth]{\FIG/dm_heatmap_Ridge_vs_DCC.png}
        \end{column}
        \begin{column}{0.48\linewidth}
            \footnotesize
            \begin{itemize}
                \item Мањи RMSE у свим приказаним експериментима: 0.066 наспрам 0.214, приближно три пута мање
                \item DM статистике са нултим лагом су описне, а не формални доказ значајности: преклапајући прозори аутокорелишу разлике губитка
                \item DCC-GARCH процењује \emph{тренутну} условну корелацију, а циљ је изглађена склизна корелација, што уводи систематско кашњење
                \item Традиционалне серије су такође прилагођене Bitcoin-овом седмодневном календару: око 30\% посматрања за GARCH компоненте су викенд-нуле
            \end{itemize}
            \normalsize
        \end{column}
    \end{columns}
\end{frame}

\begin{frame}{Слој сигнала за инвеститора}
    \begin{columns}
        \begin{column}{0.48\linewidth}
            \small
            \begin{itemize}
                \item Стресни дан: принос наредног дана испод $-0.75\,\sigma_t$; праг упозорења $\geq 0.55$
                \item Logit у просеку: BA 0.511, F1\textsubscript{пад} 0.214, AUC 0.513
                \item Најбољи случај: BTC--S\&P\,500, $w = 30$, BA 0.5308, F1 0.2430; максимум међу 20 конфигурација
                \item Random Forest дегенерише (F1 0.042): вероватноће нису калибрисане и ретко прелазе праг
            \end{itemize}
            \normalsize
        \end{column}
        \begin{column}{0.52\linewidth}
            \includegraphics[width=\linewidth]{\FIG/signal_corr_BTC-USD_^GSPC_w30.png}
        \end{column}
    \end{columns}
    \medskip
    {\tiny Регуларизовани Logit се боље генерализује од сложеног ансамбла: слаб сигнал и корелисани предиктори.}
\end{frame}

\begin{frame}{Да ли сигнал има практичну вредност?}
    \small
    \begin{center}
    \begin{tabular}{lrrrr}
        \toprule
        Пар актива & Означено (bp) & Без сигнала (bp) & Разлика & Стопа сигнала (\%) \\
        \midrule
        BTC--S\&P 500  & 4.08 & 5.05 & $+0.96$ & 33.1 \\
        BTC--NASDAQ    & 4.70 & 6.64 & $+1.94$ & 32.0 \\
        BTC--злато     & 7.01 & 7.10 & $+0.08$ & 20.1 \\
        BTC--сребро    & 10.23 & 10.13 & $-0.10$ & 26.4 \\
        BTC--USD индекс & 0.57 & 0.89 & $+0.32$ & 31.4 \\
        \bottomrule
    \end{tabular}
    \end{center}
    \medskip
    \begin{itemize}
        \item Код два акцијска пара означени дани су у просеку лошији од дана без сигнала, што је и сврха сигнала
        \item Означени приноси су и даље позитивни: узорак је растуће тржиште, па сигнал смањује изложеност релативно лошијим данима, а не предвиђа губитке
        \item Резултати су пре трансакционих трошкова, уз промену позиције приближно сваког трећег дана
    \end{itemize}
\end{frame}

% ─────────────────────────────────────────────────────────────────────────────
\section{Закључци}

\begin{frame}{Шта рад показује}
    \begin{enumerate}
        \item Временски променљива међутржишна зависност Bitcoin-а и традиционалних актива \textbf{јесте} прогнозива ван узорка: шест парова, четири прозора и строг walk-forward
        \item Прогнозивост потиче из постојаности самог циља: један лаг зависности већ даје $R^2 = 0.885$
        \item Машинско учење има мањи RMSE од DCC-GARCH модела без цурења информација у сваком приказаном експерименту, али је поређење описно; AR(1) није доследно надмашен
        \item Слој сигнала је скроман додатни индикатор ризика, најјачи за широку изложеност акцијама, а не самостално правило темпирања
    \end{enumerate}
\end{frame}

\begin{frame}{Шта рад \emph{не} тврди}
    \begin{itemize}
        \item Bitcoin не предвиђа \emph{приносе} акција. Делимично предвиђа променљиву \emph{зависност} Bitcoin-а и акција, што је ужи и прецизнији исказ
        \item Сигнал процењује вероватноћу повећаног ризика пада; не говори о смеру, величини ни времену кретања у означеном периоду
        \item Најбоља конфигурација сигнала изабрана је међу двадесет, без корекције за вишеструка поређења, па је индикативна, а не потврђена хипотеза
        \item Рад не показује да је DCC-GARCH лош модел за оно за шта је конструисан
    \end{itemize}
\end{frame}

\begin{frame}{Ограничења и будући рад}
    \begin{itemize}
        \item Дневна учесталост и хоризонт од једног корака: ауторегресивна предност би требало да слаби са дужим хоризонтом
        \item Pearson-ова склизна корелација је једина мера зависности; мера осетљива на репове тестирала би иста питања другачије
        \item Једна основна криптовалута
        \item Нема калибрације вероватноћа у слоју сигнала ни портфолио backtest-а са трансакционим трошковима
        \item Потврдни рад треба да користи угнежђено walk-forward подешавање и HAC спецификацију која моделира преклапање склизних прозора
    \end{itemize}
    \bigskip
    {\small Свако ограничење је природан правац наставка рада. Ниједно не поништава резултате у оквиру онога што је заиста тестирано.}
\end{frame}

% ─────────────────────────────────────────────────────────────────────────────
\section*{Питања}
{\setbeamertemplate{footline}{}
\begin{frame}[noframenumbering]{\secname}
    \centering
    \vfill
    
\includegraphics[width=0.34\linewidth]{slike/qr/master_thesis_site_qr.png}\\[5pt]
    {\small Интерактивна презентација рада}
    \vfill
\end{frame}}

\end{document}
